% 编译（在仓库根目录运行两遍）：
% xelatex -halt-on-error -output-directory=output/pdf assignments/assignment8/assignment8-sol.tex
% 中文字体采用 TeX Live 自带的 Fandol 字体。
\documentclass[11pt,a4paper,fontset=fandol]{ctexart}
\usepackage[margin=2.05cm]{geometry}
\usepackage{amsmath,amssymb,mathtools}
\usepackage{fancyhdr,needspace}
\usepackage[hidelinks]{hyperref}
\hypersetup{pdftitle={线性代数第八次作业},pdfsubject={MATH1205：分块矩阵、可逆性与乘积的秩}}
\setlength{\parindent}{0pt}
\setlength{\parskip}{0.2em}
\setlength{\headheight}{15pt}
\renewcommand{\arraystretch}{1.08}
\linespread{1.12}
\pagestyle{fancy}
\fancyhf{}
\fancyhead[L]{MATH1205\quad 线性代数}
\fancyhead[R]{第八次作业}
\fancyfoot[C]{\thepage}
\renewcommand{\headrulewidth}{0.3pt}
\ctexset{section={format=\large\bfseries,beforeskip=1.4ex,afterskip=0.6ex},
subsection={format=\normalsize\bfseries,beforeskip=1ex,afterskip=0.5ex}}
\newcommand{\rk}{\operatorname{r}}
\newcommand{\ans}[1]{\boxed{#1}}
\newcommand{\problem}[1]{\Needspace{9\baselineskip}\section*{#1}}

\begin{document}
\setlength{\abovedisplayskip}{7pt plus 2pt minus 2pt}
\setlength{\belowdisplayskip}{7pt plus 2pt minus 2pt}
\setlength{\abovedisplayshortskip}{4pt plus 1pt minus 1pt}
\setlength{\belowdisplayshortskip}{4pt plus 1pt minus 1pt}
\begin{center}
{\LARGE\bfseries 线性代数第八次作业}\par
\vspace{0.3em}
{\small MATH1205}
\end{center}
\vspace{-0.4em}
题号沿用原题。$E_k$ 为 $k$ 阶单位矩阵，$O$ 为相应尺寸的零矩阵，
$A^*$ 为 $A$ 的伴随矩阵，$\rk(A)$ 表示矩阵 $A$ 的秩。

\problem{第 51 题：分块矩阵的乘积与可逆性}
已知 $A$ 为 $n$ 阶可逆矩阵，$\alpha$ 为 $n\times1$ 列矩阵，$b$ 为常数，且
\[
P=\begin{pmatrix}E_n&O\\-\alpha^T A^*&|A|\end{pmatrix},
\qquad
Q=\begin{pmatrix}A&\alpha\\\alpha^T&b\end{pmatrix}.
\]
\textbf{(1) 计算并化简 $PQ$。}
利用伴随矩阵的恒等式 $A^*A=|A|E_n$，得
\[
\begin{aligned}
PQ
&=\begin{pmatrix}
A&\alpha\\
-\alpha^T A^*A+|A|\alpha^T&-\alpha^T A^*\alpha+b|A|
\end{pmatrix}\\[2pt]
&=\ans{\begin{pmatrix}
A&\alpha\\
O&b|A|-\alpha^T A^*\alpha
\end{pmatrix}}.
\end{aligned}
\]
又因 $A^*=|A|A^{-1}$，右下角也可写为
$|A|\bigl(b-\alpha^T A^{-1}\alpha\bigr)$。

\textbf{(2) 证明 $Q$ 可逆的充分必要条件。}
矩阵 $P$ 为分块下三角矩阵，故
\[
|P|=|E_n|\,|A|=|A|\ne0.
\]
由第 (1) 问及分块上三角矩阵的行列式公式，
\[
|P|\,|Q|=|PQ|
=|A|^2\bigl(b-\alpha^T A^{-1}\alpha\bigr).
\]
约去非零因子 $|A|$，得到
\[
|Q|=|A|\bigl(b-\alpha^T A^{-1}\alpha\bigr).
\]
因此
\[
\ans{Q\text{ 可逆}\iff |Q|\ne0
\iff \alpha^T A^{-1}\alpha\ne b}.
\]

\problem{第 68 题：由 $E-AB$ 可逆推出 $E-BA$ 可逆}
设 $C=(E-AB)^{-1}$，并令 $D=E+BCA$。由
\[
(E-BA)B=B(E-AB)
\]
可得
\[
\begin{aligned}
(E-BA)D
&=(E-BA)(E+BCA)\\
&=E-BA+B(E-AB)CA\\
&=E-BA+BA=E.
\end{aligned}
\]
于是 $|E-BA|\,|D|=1$，从而 $|E-BA|\ne0$，即 $E-BA$ 可逆。
同时得到
\[
\ans{(E-BA)^{-1}=E+B(E-AB)^{-1}A}.
\]

\newpage
\problem{第 69 题：证明 $|\lambda E-AB|=|\lambda E-BA|$}
其中 $A,B$ 均为 $n$ 阶方阵，$\lambda\ne0$。
构造 $2n$ 阶分块矩阵
\[
M=\begin{pmatrix}\lambda E-AB&A\\O&\lambda E\end{pmatrix}.
\]
先将第二块列右乘 $B$ 加到第一块列，再将第一块行左乘 $-B$ 加到第二块行：
\[
\begin{pmatrix}\lambda E-AB&A\\O&\lambda E\end{pmatrix}
\xrightarrow{C_1\leftarrow C_1+C_2B}
\begin{pmatrix}\lambda E&A\\\lambda B&\lambda E\end{pmatrix}
\xrightarrow{R_2\leftarrow R_2-BR_1}
\begin{pmatrix}\lambda E&A\\O&\lambda E-BA\end{pmatrix}.
\]
这里 $R_i,C_j$ 分别表示块行、块列。两次变换均为第三类分块初等变换，
不改变行列式，因此
\[
|\lambda E-AB|\,|\lambda E|
=|\lambda E|\,|\lambda E-BA|.
\]
由于 $|\lambda E|=\lambda^n\ne0$，约去该因子，得
\[
\ans{|\lambda E-AB|=|\lambda E-BA|}.
\]

\problem{第 70 题：矩阵乘积的秩不等式}
设 $A$ 为 $m\times n$ 矩阵，$B$ 为 $n\times s$ 矩阵。证明
\[
\rk(A)+\rk(B)-n\le\rk(AB)\le\min\{\rk(A),\rk(B)\}.
\]
\textbf{(1) 上界。}
对分块矩阵作初等变换：
\[
\begin{pmatrix}A&AB\end{pmatrix}
\xrightarrow{C_2\leftarrow C_2-C_1B}
\begin{pmatrix}A&O\end{pmatrix},
\qquad
\begin{pmatrix}B\\AB\end{pmatrix}
\xrightarrow{R_2\leftarrow R_2-AR_1}
\begin{pmatrix}B\\O\end{pmatrix}.
\]
初等变换保持秩，且子矩阵的秩不超过原矩阵的秩，所以
\[
\rk(AB)\le\rk\begin{pmatrix}A&AB\end{pmatrix}=\rk(A),
\qquad
\rk(AB)\le\rk\begin{pmatrix}B\\AB\end{pmatrix}=\rk(B).
\]
故 $\rk(AB)\le\min\{\rk(A),\rk(B)\}$。

\textbf{(2) 下界。}
记 $r=\rk(A)$。由矩阵的等价标准形，存在 $m$ 阶、$n$ 阶可逆矩阵 $U,V$，使
\[
UAV=\begin{pmatrix}E_r&O\\O&O\end{pmatrix}.
\]
将 $V^{-1}B$ 按前 $r$ 行和其余 $n-r$ 行分块，记作
\[
V^{-1}B=\begin{pmatrix}B_1\\B_2\end{pmatrix},
\qquad B_1\in\mathbb{F}^{r\times s},\quad
B_2\in\mathbb{F}^{(n-r)\times s}.
\]
于是
\[
UAB=(UAV)(V^{-1}B)
=\begin{pmatrix}E_r&O\\O&O\end{pmatrix}
\begin{pmatrix}B_1\\B_2\end{pmatrix}
=\begin{pmatrix}B_1\\O\end{pmatrix},
\qquad \rk(AB)=\rk(B_1).
\]
只在 $B_1$ 的行之间作初等行变换，将其化为阶梯形。这一部分恰有
$\rk(B_1)$ 个非零行，而 $B_2$ 部分至多有 $n-r$ 个非零行。因此
\[
\rk(B)=\rk(V^{-1}B)
\le\rk(B_1)+(n-r)
=\rk(AB)+n-\rk(A).
\]
移项并结合上界，即得
\[
\ans{\rk(A)+\rk(B)-n\le\rk(AB)
\le\min\{\rk(A),\rk(B)\}}.
\]
上述分块允许零行或零列；$r=0$ 或 $r=n$ 时结论同样成立。

\end{document}
